% hierarchy.tex - Hierarchy and tree diagram macros for Paperforge
% Requires: _colors.tex, _styles.tex loaded first
%
% Macros provided:
%   \pyramid{top}{middle}{bottom}
%   \pyramidfour{t1}{t2}{t3}{t4}
%   \orgtree{root}{child1}{child2}
%   \orgtreethree{root}{c1}{c2}{c3}
%   \treediagram - manual tree construction
%
% Design: Pyramids use gradient tiers, trees use hierarchical spacing

% ============================================================================
% THREE-TIER PYRAMID
% ============================================================================
% Usage: \pyramid{Top Tier}{Middle Tier}{Bottom Tier}
\newcommand{\pyramid}[3]{%
\begin{center}
\begin{tikzpicture}
    % Top tier (smallest)
    \node[trapezium, trapezium angle=70, draw=diagramTier1, fill=diagramTier1!20,
          minimum width=2.5cm, minimum height=0.9cm, align=center, font=\small\bfseries,
          text=diagramTier1] (t1) at (0,2) {#1};

    % Middle tier
    \node[trapezium, trapezium angle=70, draw=diagramTier2, fill=diagramTier2!20,
          minimum width=5cm, minimum height=0.9cm, align=center, font=\small\bfseries,
          text=diagramTier2] (t2) at (0,1) {#2};

    % Bottom tier (largest)
    \node[trapezium, trapezium angle=70, draw=diagramTier3, fill=diagramTier3!20,
          minimum width=7.5cm, minimum height=0.9cm, align=center, font=\small\bfseries,
          text=diagramTier3] (t3) at (0,0) {#3};
\end{tikzpicture}
\end{center}
}

% ============================================================================
% FOUR-TIER PYRAMID
% ============================================================================
% Usage: \pyramidfour{Tier 1}{Tier 2}{Tier 3}{Tier 4}
\newcommand{\pyramidfour}[4]{%
\begin{center}
\begin{tikzpicture}
    % Top tier
    \node[trapezium, trapezium angle=70, draw=diagramTier1, fill=diagramTier1!20,
          minimum width=2cm, minimum height=0.8cm, align=center, font=\small\bfseries,
          text=diagramTier1] (t1) at (0,2.7) {#1};

    % Second tier
    \node[trapezium, trapezium angle=70, draw=diagramTier2, fill=diagramTier2!20,
          minimum width=4cm, minimum height=0.8cm, align=center, font=\small\bfseries,
          text=diagramTier2] (t2) at (0,1.8) {#2};

    % Third tier
    \node[trapezium, trapezium angle=70, draw=diagramTier3, fill=diagramTier3!20,
          minimum width=6cm, minimum height=0.8cm, align=center, font=\small\bfseries,
          text=diagramTier3] (t3) at (0,0.9) {#3};

    % Bottom tier
    \node[trapezium, trapezium angle=70, draw=diagramTier4, fill=diagramTier4!50,
          minimum width=8cm, minimum height=0.8cm, align=center, font=\small\bfseries,
          text=diagramTextDark] (t4) at (0,0) {#4};
\end{tikzpicture}
\end{center}
}

% ============================================================================
% INVERTED PYRAMID (funnel)
% ============================================================================
% Usage: \funnel{Wide Top}{Middle}{Narrow Bottom}
\newcommand{\funnel}[3]{%
\begin{center}
\begin{tikzpicture}
    % Top (widest)
    \node[trapezium, trapezium angle=110, draw=diagramTier1, fill=diagramTier1!20,
          minimum width=7.5cm, minimum height=0.9cm, align=center, font=\small\bfseries,
          text=diagramTier1] (t1) at (0,2) {#1};

    % Middle
    \node[trapezium, trapezium angle=110, draw=diagramTier2, fill=diagramTier2!20,
          minimum width=5cm, minimum height=0.9cm, align=center, font=\small\bfseries,
          text=diagramTier2] (t2) at (0,1) {#2};

    % Bottom (narrowest)
    \node[trapezium, trapezium angle=110, draw=diagramTier3, fill=diagramTier3!20,
          minimum width=2.5cm, minimum height=0.9cm, align=center, font=\small\bfseries,
          text=diagramTier3] (t3) at (0,0) {#3};

    % Arrow indicating flow
    \draw[diagram arrow, color=diagramNeutral] (0, -0.6) -- (0, -1);
\end{tikzpicture}
\end{center}
}

% ============================================================================
% ORG CHART TREE (2 children)
% ============================================================================
% Usage: \orgtree{Root}{Child 1}{Child 2}
\newcommand{\orgtree}[3]{%
\begin{center}
\begin{tikzpicture}[
    level distance=1.8cm,
    sibling distance=4cm,
    edge from parent/.style={diagram arrow}
]
    \node[diagram tree node, fill=diagramPrimary!15, draw=diagramPrimary, font=\small\bfseries] {#1}
        child {node[diagram tree node] {#2}}
        child {node[diagram tree node] {#3}};
\end{tikzpicture}
\end{center}
}

% ============================================================================
% ORG CHART TREE (3 children)
% ============================================================================
% Usage: \orgtreethree{Root}{Child 1}{Child 2}{Child 3}
\newcommand{\orgtreethree}[4]{%
\begin{center}
\begin{tikzpicture}[
    level distance=1.8cm,
    sibling distance=3.5cm,
    edge from parent/.style={diagram arrow}
]
    \node[diagram tree node, fill=diagramPrimary!15, draw=diagramPrimary, font=\small\bfseries] {#1}
        child {node[diagram tree node] {#2}}
        child {node[diagram tree node] {#3}}
        child {node[diagram tree node] {#4}};
\end{tikzpicture}
\end{center}
}

% ============================================================================
% TWO-LEVEL ORG CHART (Root → 2 → 4)
% ============================================================================
% Usage: \orgcharttwolevel{Root}{Manager1}{Manager2}{Staff1}{Staff2}{Staff3}{Staff4}
\newcommand{\orgcharttwolevel}[7]{%
\begin{center}
\begin{tikzpicture}[
    level 1/.style={sibling distance=5.5cm, level distance=1.5cm},
    level 2/.style={sibling distance=2.5cm, level distance=1.5cm},
    edge from parent/.style={diagram arrow}
]
    \node[diagram tree node, fill=diagramPrimary!15, draw=diagramPrimary, font=\small\bfseries] {#1}
        child {node[diagram tree node, fill=diagramSecondary!15] {#2}
            child {node[diagram tree node] {#4}}
            child {node[diagram tree node] {#5}}
        }
        child {node[diagram tree node, fill=diagramSecondary!15] {#3}
            child {node[diagram tree node] {#6}}
            child {node[diagram tree node] {#7}}
        };
\end{tikzpicture}
\end{center}
}

% ============================================================================
% SIMPLE HIERARCHY LIST (vertical)
% ============================================================================
% Usage: \hierarchylist{Level 1}{Level 2}{Level 3}
\newcommand{\hierarchylist}[3]{%
\begin{center}
\begin{tikzpicture}
    % Level 1 (top)
    \node[diagram box primary, minimum width=6cm] (l1) at (0, 0) {#1};

    % Level 2
    \node[diagram box, minimum width=6cm, below=0.8cm of l1] (l2) {#2};

    % Level 3
    \node[diagram box, minimum width=6cm, below=0.8cm of l2] (l3) {#3};

    % Connecting lines
    \draw[diagram arrow] (l1) -- (l2);
    \draw[diagram arrow] (l2) -- (l3);
\end{tikzpicture}
\end{center}
}

% ============================================================================
% DARK THEME PYRAMID
% ============================================================================
% Usage: \pyramiddark{Top}{Middle}{Bottom}
\newcommand{\pyramiddark}[3]{%
\begin{center}
\begin{tikzpicture}
    % Top tier
    \node[trapezium, trapezium angle=70, draw=diagramDarkBorder, fill=diagramDarkPrimary!30,
          minimum width=2.5cm, minimum height=0.9cm, align=center, font=\small\bfseries,
          text=diagramDarkText] (t1) at (0,2) {#1};

    % Middle tier
    \node[trapezium, trapezium angle=70, draw=diagramDarkBorder, fill=diagramDarkSurface,
          minimum width=5cm, minimum height=0.9cm, align=center, font=\small\bfseries,
          text=diagramDarkText] (t2) at (0,1) {#2};

    % Bottom tier
    \node[trapezium, trapezium angle=70, draw=diagramDarkBorder, fill=diagramDarkBg,
          minimum width=7.5cm, minimum height=0.9cm, align=center, font=\small\bfseries,
          text=diagramDarkText] (t3) at (0,0) {#3};
\end{tikzpicture}
\end{center}
}
